problem:  
Let \( X \) and \( X^\vee \) be compact Calabi–Yau threefolds over \( \mathbb{C} \) satisfying homological mirror symmetry, i.e. \( D^b\!\operatorname{Coh}(X)\simeq D^\pi\!\operatorname{Fuk}(X^\vee) \).  
For each, consider the *derived moduli stack of simple perfect complexes*  
\[
\mathfrak{M}(X)\,,\qquad  \mathfrak{M}(X^\vee),
\]
whose truncations parametrize isomorphism classes of simple objects of finite amplitude.  
By Pantev–Toën–Vaquié–Vezzosi, these stacks carry canonical \((-1)\)-shifted symplectic forms \(\omega_X, \omega_{X^\vee}\).  For a Fourier–Mukai kernel \(P\in D^b\!\operatorname{Coh}(X\times X^\vee)\), one obtains an induced morphism of derived stacks \( \Phi_P:\mathfrak{M}(X)\to \mathfrak{M}(X^\vee)\).

**Problem:**  
Suppose homological mirror symmetry holds for \(X, X^\vee\).  Does there exist a kernel \(P\) such that the induced morphism \(\Phi_P\) is a *Lagrangian morphism* between their \((-1)\)-shifted symplectic structures, i.e.  
\[
\Phi_P^*\omega_{X^\vee} = 0 \quad \text{in } H^{-1}(\Omega_{\mathfrak{M}(X)})\,,  
\]
and the relative obstruction theory of \(\Phi_P\) realizes the mirror correspondence of deformation–obstruction complexes?  
Equivalently, can the homological mirror functor be represented by a Fourier–Mukai kernel whose induced morphism between moduli of objects is Lagrangian in the sense of shifted symplectic geometry?

Why is it a "good" problem:  
This version removes undefined notions and builds entirely on rigorously available structures—derived moduli stacks of simple objects and their canonical shifted symplectic forms.  The question asks whether the *categorical phenomenon of homological mirror symmetry* admits a *geometric manifestation* inside shifted symplectic geometry, namely as a Lagrangian morphism between the moduli of objects on mirror Calabi–Yau threefolds.  The impact would be major: it would establish a direct geometric link between PTVV’s theory of shifted symplectic structures and Kontsevich’s categorical mirror symmetry.  The formulation is simple and precise yet open; it bridges algebraic geometry, category theory, and derived symplectic geometry, representing a deep and coherent research‑level problem.